\chapter{LASIK}

\section*{Introduction \& Clinical Indications}
LASIK, or laser-assisted in situ keratomileusis, is a popular elective refractive surgical procedure aimed at correcting common vision disorders, including myopia (nearsightedness), hyperopia (farsightedness), and astigmatism. This innovative technique utilizes a specialized excimer laser to reshape the corneal stroma, the transparent front part of the eye, after creating a hinged flap in the cornea. By modifying the curvature of the cornea, LASIK alters the eye's focusing power, allowing light to focus accurately on the retina, which can significantly reduce or eliminate the need for corrective lenses. The procedure is performed on an outpatient basis and is known for its rapid recovery and high success rates, making it a transformative option for individuals seeking visual independence and improved quality of life.

\begin{itemize}
  \item Laser Eye Surgery
  \item Laser Vision Correction
  \item Corneal Refractive Surgery
  \item Excimer Laser Keratomileusis
  \item Refractive Laser Surgery
  \item Femtosecond LASIK (when femtosecond laser is used for flap creation)
\end{itemize}

The fundamental principle of LASIK involves modifying the anterior curvature of the corneal stroma to alter its refractive power, thereby correcting the eye's ability to focus light onto the retina. The procedure begins with the creation of a thin, hinged corneal flap, which is then gently lifted to expose the underlying corneal stroma. This flap can be created either mechanically using a microkeratome blade or, more commonly, with a femtosecond laser, which generates precise micro-bubbles within the corneal tissue to separate the layers. The rationale for creating a flap is to preserve the corneal epithelium and Bowman's layer, which allows for faster visual recovery, reduced postoperative pain, and less risk of corneal haze compared to surface ablation techniques like PRK.

Once the stromal bed is exposed, an excimer laser is used to ablate (remove) microscopic amounts of corneal tissue with submicron precision. For myopia, tissue is removed from the central cornea to flatten its curvature, effectively reducing the eye's focusing power. For hyperopia, tissue is removed from the periphery to steepen the central curvature, thereby increasing the eye's focusing power. For astigmatism, tissue is removed unevenly in a specific pattern to make the cornea more spherical and eliminate the irregular focusing. The laser's pulses break molecular bonds in the corneal tissue through photoablation, without causing significant thermal damage to surrounding cells. After the desired reshaping is achieved, the corneal flap is carefully repositioned over the ablated stroma, where it adheres naturally without sutures, acting as a biological bandage. The technical goal is to achieve emmetropia (perfect focus) or a target refraction that provides optimal uncorrected visual acuity, tailored to the patient's visual needs and lifestyle.

The journey towards modern LASIK began with the development of the excimer laser in the 1970s, specifically its application to biological tissue by Rangaswamy Srinivasan and Stephen Trokel in the early 1980s, demonstrating its ability to precisely ablate corneal tissue without thermal damage. This groundbreaking discovery led to the first photorefractive keratectomy (PRK) procedures, where the corneal surface epithelium was removed before laser ablation, pioneered by Theo Seiler and Marguerite McDonald in the late 1980s. PRK offered effective refractive correction but was associated with significant postoperative pain and slower visual recovery.

The concept of creating a corneal flap to protect the ablated stroma and facilitate faster healing was explored much earlier by Jose Barraquer in the 1950s with his keratomileusis technique, which involved freezing and reshaping corneal tissue. However, it was Ioannis Pallikaris in Greece who, in 1990, combined the excimer laser ablation with a hinged corneal flap created by a microkeratome, coining the term LASIK. Lucio Buratto also contributed significantly to the early development and refinement of the LASIK technique. The introduction of wavefront-guided LASIK in the early 2000s further enhanced precision by correcting higher-order aberrations, leading to improved visual quality. The advent of femtosecond lasers for flap creation around the same time significantly improved flap safety and predictability, reducing microkeratome-related complications and solidifying LASIK's position as the most common elective refractive procedure globally.

LASIK is indicated for individuals seeking to reduce or eliminate their dependence on glasses or contact lenses, provided they meet specific criteria for ocular health and refractive stability.

\begin{itemize}
  \item To correct myopia (nearsightedness) typically ranging from -1.00 to -10.00 diopters, reducing dependence on glasses or contact lenses.
  \item To correct hyperopia (farsightedness) typically ranging from +1.00 to +5.00 diopters, improving distance and often near vision without corrective eyewear.
  \item To correct astigmatism, which causes blurred vision at all distances due to an irregularly shaped cornea, typically up to -5.00 or -6.00 diopters.
  \item For individuals seeking freedom from daily use of glasses or contact lenses for lifestyle, occupational (e.g., military, sports), or cosmetic reasons.
  \item When contact lens intolerance, chronic dry eye symptoms, or recurrent infections are exacerbated by contact lens wear.
  \item For patients with stable refractive error for at least one year, indicating mature and consistent vision, typically aged 18 years or older.
  \item For individuals who meet specific corneal health criteria, including adequate corneal thickness and absence of corneal disease or ectasia.
  \item To improve uncorrected visual acuity and overall quality of life for suitable candidates who understand the risks and benefits.
\end{itemize}

LASIK is overwhelmingly considered a primary elective refractive procedure for suitable candidates. For individuals who meet the stringent criteria for corneal health, refractive stability, and overall ocular health, it is often the first-line surgical option chosen to correct refractive errors. Its rapid visual recovery, minimal discomfort, and high predictability make it a preferred choice over other refractive surgeries for a broad range of patients.

While LASIK is primarily a first-line procedure, in some rare instances, an "enhancement" LASIK (or PRK) may be performed years later to correct residual refractive error or regression, effectively acting as a secondary procedure. However, for initial refractive correction, it stands as a primary choice, offering a high degree of predictability and rapid visual recovery compared to other refractive surgical options like phakic intraocular lenses (IOLs) or refractive lens exchange (RLE). Phakic IOLs are generally reserved for cases outside the LASIK treatment range (e.g., very high myopia) or with specific contraindications, while RLE is typically considered for older patients with presbyopia or early cataracts.

\section*{Relevant Surgical Anatomy}
The cornea is the primary refractive surface of the eye, accounting for approximately two-thirds of its total refractive power. It is a transparent, avascular, dome-shaped structure composed of five distinct layers: the epithelium, Bowman's layer, the corneal stroma, Descemet's membrane, and the endothelium. The corneal stroma, which constitutes about 90\% of the corneal thickness, is the target tissue for excimer laser ablation in LASIK. This layer is composed of highly organized collagen lamellae and keratocytes, which are essential for maintaining transparency and structural integrity. The average central corneal thickness is approximately 550 microns, a critical measurement for determining LASIK candidacy and ensuring adequate residual stromal bed thickness after ablation.

The cornea's nerve supply is primarily derived from the ophthalmic division of the trigeminal nerve (cranial nerve V), specifically the long ciliary nerves. These nerves penetrate the limbus (the junction between the cornea and sclera) and form a dense subepithelial and stromal plexus. Severing these nerves during flap creation is a key factor contributing to postoperative dry eye syndrome. The cornea receives its nutrients from the aqueous humor, tear film, and the limbal capillary plexus at its periphery. Surgical landmarks include the limbus, which defines the corneal-scleral boundary, and the pupil, which represents the optical center of the eye and guides the centration of laser ablation. The visual axis, which passes through the center of the pupil and fovea, is the critical reference for achieving optimal refractive correction.

\section*{Preoperative Workup \& Preparation}
Thorough preoperative evaluation is critical for determining candidacy and ensuring optimal outcomes for LASIK. This comprehensive assessment aims to identify any contraindications, accurately measure the refractive error, and map the corneal topography.

Key components of the preoperative workup include:

\begin{itemize}
  \item \textbf{Comprehensive Eye Examination:} This includes a detailed medical and ocular history, measurement of uncorrected and best-corrected visual acuity, manifest and cycloplegic refraction (to determine the true refractive error without accommodation), intraocular pressure measurement, and a dilated fundus examination to rule out retinal pathology.
  
  \item \textbf{Corneal Diagnostics:} Essential imaging includes corneal topography and tomography (e.g., Pentacam, Orbscan) to map the corneal curvature, elevation, and detect subtle irregularities or signs of ectatic disease (like keratoconus or forme fruste keratoconus). Corneal pachymetry (thickness measurement) is crucial to ensure sufficient residual stromal bed thickness (typically >250-300 microns) after ablation. Wavefront analysis may also be performed to measure higher-order aberrations, allowing for customized ablation profiles.
  
  \item \textbf{Contact Lens Discontinuation:} Patients must discontinue soft contact lens wear for at least 1-2 weeks and rigid gas permeable (RGP) lenses for 3-4 weeks (or longer, depending on corneal stability) prior to the preoperative evaluation and surgery. This allows the cornea to return to its natural shape, preventing inaccurate measurements and potential complications.
  
  \item \textbf{Medication Review:} Patients should inform their surgeon of all systemic and topical medications, as some, like isotretinoin, can cause severe dry eye and may need to be discontinued or adjusted.
  
  \item \textbf{Informed Consent:} A detailed discussion of the procedure, potential risks, benefits, alternatives (e.g., PRK, phakic IOLs), and expected outcomes is conducted. The patient must fully understand and provide informed consent.
  
  \item \textbf{Prophylactic Antibiotics:} Patients may be prescribed topical antibiotic eye drops to start a day or two before surgery to minimize the risk of bacterial infection.
\end{itemize}

No specific dietary restrictions or NPO (nil per os) status are typically required, as only topical anesthesia is used, and the procedure is brief. Patients are advised to arrange for transportation home after the procedure.

\textbf{Situations requiring treatment, delay, or an alternative refractive option:}

\begin{itemize}
  \item Keratoconus or other corneal ectatic disorders (e.g., pellucid marginal degeneration), as LASIK can significantly worsen corneal instability and lead to progressive vision loss.
  
  \item Insufficient corneal thickness for the planned ablation, which would result in an unsafe residual stromal bed (typically less than 250-300 microns).
  
  \item Unstable refractive error, indicating ongoing changes in vision that would lead to suboptimal long-term results and potential regression.
  
  \item Active ocular infection (e.g., herpes simplex keratitis, bacterial conjunctivitis) or inflammation (e.g., uveitis).
  
  \item Severe dry eye syndrome that is unresponsive to aggressive medical treatment, as LASIK can exacerbate symptoms significantly.
  
  \item Autoimmune diseases (e.g., rheumatoid arthritis, lupus, Sj\"ogren's syndrome) or other systemic conditions that impair wound healing (e.g., uncontrolled diabetes).
  
  \item Pregnancy or nursing, due to hormonal fluctuations affecting corneal hydration and refraction, and potential medication risks to the fetus/infant.
  
  \item Uncontrolled glaucoma or advanced cataracts that would limit the potential visual benefit of LASIK.
  
  \item History of previous corneal surgery (e.g., radial keratotomy, corneal transplant) that may compromise corneal integrity or healing.
\end{itemize}

\textbf{Relative Precautions:}

\begin{itemize}
  \item Moderate dry eye syndrome, which requires aggressive preoperative and postoperative management, including punctal plugs or prescription eye drops.
  
  \item Large pupil size (greater than 6.5-7.0 mm in dim light), which may increase the risk of night vision disturbances like glare, halos, and starbursts due to the optical zone being smaller than the pupil.
  
  \item Certain systemic medications (e.g., corticosteroids, amiodarone, chloroquine) that can affect corneal healing, cause corneal deposits, or alter vision.
  
  \item High refractive errors (e.g., myopia >-8.00 D, hyperopia >+4.00 D), which may increase the risk of regression, ectasia, or suboptimal visual outcomes.
  
  \item Age under 18-21 years, as refractive error may not yet be stable.
  
  \item Monocular vision (vision in only one eye), due to the higher risk associated with the only functional eye.
  
  \item Blepharitis or meibomian gland dysfunction, which should be treated preoperatively to reduce infection risk and improve dry eye symptoms.
\end{itemize}

\section*{Surgical Technique \& Procedure}
LASIK is performed as an outpatient procedure, typically in a dedicated laser suite, under topical anesthesia. The patient is positioned supine on a treatment bed, and the head is stabilized. The eye to be treated is prepped with antiseptic drops (e.g., povidone-iodine), and an eyelid speculum is inserted to keep the eye open and prevent blinking.

The procedure involves several key steps:

\begin{enumerate}
  \item \textbf{Anesthesia and Positioning:} Topical anesthetic drops (e.g., proparacaine, tetracaine) are applied to numb the eye. The patient is positioned comfortably, and the head is aligned under the excimer laser system.
  
  \item \textbf{Eye Preparation and Suction Ring Application:} The eye is cleaned, and a suction ring is carefully applied to the globe. This ring stabilizes the eye and elevates intraocular pressure, which temporarily darkens vision and ensures a stable platform for flap creation.
  
  \item \textbf{Corneal Flap Creation:}
      \begin{itemize}
        \item \textit{Microkeratome Method:} A mechanical microkeratome, a precision oscillating blade, is guided across the cornea to create a thin, hinged flap (typically 100-160 microns thick).
        
        \item \textit{Femtosecond Laser Method (bladeless LASIK):} A femtosecond laser delivers rapid, ultra-short pulses of light to create a series of tiny, overlapping gas bubbles within the corneal stroma, precisely separating a layer of tissue to form the flap. This method offers greater precision and reduced risk of certain flap complications. The flap is typically hinged superiorly or nasally.
      \end{itemize}
  
  \item \textbf{Flap Lifting and Stromal Exposure:} Once the flap is created, the suction ring is removed, and the surgeon gently lifts and folds back the corneal flap, exposing the underlying stromal bed.
  
  \item \textbf{Excimer Laser Ablation:} The patient is instructed to fixate on a target light, and an integrated eye-tracking system ensures the laser pulses are delivered accurately despite minor eye movements. The excimer laser then ablates the corneal tissue according to the patient's specific refractive error profile, which was pre-programmed from the preoperative measurements. This photoablation process usually takes less than a minute per eye.
  
  \item \textbf{Flap Repositioning and Adherence:} After the laser ablation is complete, the corneal flap is carefully repositioned onto the stromal bed. The surgeon ensures proper alignment and gently irrigates the interface to remove any debris or epithelial cells. The flap adheres naturally within minutes due to osmotic forces and the smooth interface, without the need for sutures.
  
  \item \textbf{Post-Procedure Care:} Both eyes are typically treated in the same session, often sequentially. Protective shields or goggles are placed over the eyes, and the patient is given detailed instructions for postoperative care, including the use of eye drops.
\end{enumerate}

\section*{Complications \& Risks}
While LASIK is generally safe and effective, like any surgical procedure, it carries potential risks and complications. These can be broadly categorized into general surgical risks and procedure-specific risks.

\textbf{General Surgical Complications:}

\begin{itemize}
  \item \textbf{Infection:} Although rare (incidence <0.1\%), bacterial or fungal infection of the cornea (keratitis) can occur and may lead to severe vision loss if not promptly treated.
  
  \item \textbf{Inflammation:} Postoperative inflammation, while usually mild and controlled with steroids, can sometimes be more severe.
  
  \item \textbf{Bleeding:} Minimal for LASIK, typically limited to subconjunctival hemorrhage from the suction ring, which resolves spontaneously.
  
  \item \textbf{Anesthesia Complications:} Minimal for topical anesthesia, but rare allergic reactions or systemic effects are possible.
\end{itemize}

\textbf{Procedure-Specific Complications:}

\begin{itemize}
  \item \textbf{Dry Eye Syndrome:} This is the most common complication, often temporary but can be chronic. It results from severed corneal nerves during flap creation, reducing tear production and sensation, and can lead to discomfort, foreign body sensation, and fluctuating vision.
  
  \item \textbf{Glare, Halos, and Starbursts:} These night vision disturbances can occur, especially in patients with large pupils or high refractive errors, due to the optical zone being smaller than the pupil in dim light or due to induced higher-order aberrations.
  
  \item \textbf{Flap Complications:} These are rare with modern femtosecond laser techniques but can include:
    \begin{itemize}
      \item Incomplete flap or buttonhole (more common with microkeratome).
      \item Free cap (flap completely detached, more common with microkeratome).
      \item Displaced or wrinkled flap, requiring immediate repositioning.
      \item Epithelial ingrowth, where surface cells grow under the flap, potentially requiring surgical removal.
      \item Diffuse Lamellar Keratitis (DLK), also known as "sands of Sahara," an inflammatory reaction under the flap that can affect vision if severe.
    \end{itemize}
  
  \item \textbf{Undercorrection or Overcorrection:} The laser ablation may not achieve the exact target refraction, leading to residual myopia, hyperopia, or astigmatism, sometimes requiring an enhancement procedure.
  
  \item \textbf{Regression:} Over time, the cornea may partially revert to its original shape, leading to a return of some refractive error, particularly in high corrections.
  
  \item \textbf{Corneal Ectasia:} A rare but severe complication where the cornea progressively thins and bulges forward, similar to keratoconus, leading to irregular astigmatism and significant vision loss. This risk is higher in patients with undiagnosed forme fruste keratoconus or insufficient residual stromal bed.
  
  \item \textbf{Induced Astigmatism or Irregular Astigmatism:} Can result from decentered ablation, uneven healing, or flap irregularities, causing distorted vision that may not be correctable with glasses.
  
  \item \textbf{Vision Loss:} While extremely rare, severe complications such as intractable infection, progressive ectasia, or severe flap issues can lead to permanent vision impairment or even blindness.
\end{itemize}

\section*{Postoperative Care \& Recovery}
Immediately after LASIK, patients typically experience a mild burning, gritty sensation, or foreign body sensation in their eyes, which usually subsides within a few hours. Vision will be blurry but often shows significant improvement within the first 24 hours. Patients are advised to rest their eyes, avoid rubbing them, and wear protective eye shields or goggles for the first night and during sleep for about a week to prevent accidental trauma to the healing flap.

For the first 24-48 hours, patients are prescribed antibiotic eye drops to prevent infection and steroid eye drops to reduce inflammation, typically used for about one week. Lubricating eye drops are crucial and often used frequently for several weeks to months to manage postoperative dry eye symptoms. Most patients can resume light activities and return to work within 1-3 days, though activities like swimming, hot tubs, saunas, and contact sports are restricted for several weeks (typically 2-4 weeks) to prevent infection or flap displacement. Vision continues to stabilize over the first few weeks to months, with the final visual outcome typically achieved by 3-6 months post-surgery as the cornea fully heals and nerve regeneration occurs. Long-term, patients should continue to use lubricating drops as needed and attend regular eye check-ups.

During a LASIK procedure, the surgeon's primary findings relate to the integrity of the corneal flap and the precision of the laser ablation. The surgeon observes the creation of a smooth, uniform flap, its successful lifting without tears or irregularities, and the accurate delivery of the excimer laser pulses to the stromal bed. After ablation, the stromal surface should appear smooth and well-centered. Upon repositioning, the flap should seat perfectly, with minimal interface debris and no significant wrinkles or striae. These intraoperative observations are crucial for identifying and managing any immediate complications.

Unlike many other surgical procedures, LASIK does not involve the excision of tissue for histological analysis. Therefore, there is no pathology report generated from resected tissues. The "pathology" in LASIK refers to the microscopic changes induced by the laser ablation within the corneal stroma, which are designed to be precise and controlled, altering the tissue's refractive properties without causing significant cellular damage or scarring. Any adverse "pathology" would be a complication such as epithelial ingrowth or diffuse lamellar keratitis, which are clinical diagnoses rather than findings from a pathology lab.

A normal and successful postoperative course following LASIK is characterized by rapid visual recovery and minimal discomfort. Patients typically experience a significant improvement in uncorrected visual acuity within the first 24 hours, often achieving 20/20 or better vision. Any initial mild burning, grittiness, or light sensitivity usually resolves within the first day. Pain is generally minimal and well-controlled with over-the-counter analgesics. The cornea should appear clear and healthy, with the flap well-seated and no signs of inflammation, infection, or epithelial ingrowth.

Over the subsequent weeks and months, vision continues to stabilize, and any temporary symptoms such as dry eye, glare, or halos should gradually improve. The refractive error should be stable and close to emmetropia, indicating that the corneal reshaping has effectively corrected the underlying vision problem. Patients should be able to resume most normal activities, including driving, within a few days, and more strenuous activities within a few weeks. Ultimately, a normal course signifies a significant improvement in quality of life and visual freedom, allowing the patient to function without the need for glasses or contact lenses.

Postoperative care is essential for monitoring healing, preventing complications, and ensuring optimal long-term outcomes after LASIK.

\begin{itemize}
  \item \textbf{Immediate Postoperative Visits:} Typically scheduled for 1 day, 1 week, and 1 month after surgery to assess visual acuity, corneal health, flap integrity, and intraocular pressure.
  
  \item \textbf{Subsequent Follow-ups:} Further visits are usually scheduled at 3 months, 6 months, and 1 year, and then annually, to monitor long-term stability, detect any regression, and manage ongoing symptoms like dry eye.
  
  \item \textbf{Medication Management:} The use of topical antibiotic and steroid eye drops will be tapered and discontinued as directed by the surgeon, typically within the first week. Lubricating eye drops may be continued for several months or longer as needed to manage dry eye.
  
  \item \textbf{Activity Resumption:} Patients are advised on a gradual return to normal activities, with restrictions on swimming, hot tubs, saunas, and contact sports for several weeks to prevent infection or flap trauma.
  
  \item \textbf{Warning Signs:} Patients are instructed to contact their surgeon immediately if they experience sudden severe pain, significant decrease in vision, excessive redness, discharge, or new onset of severe light sensitivity, as these could indicate a serious complication.
  
  \item \textbf{Enhancement Procedures:} If a significant residual refractive error persists after 3-6 months of stable refraction and the cornea is deemed suitable, an enhancement procedure (either LASIK or PRK) may be discussed to refine the visual outcome.
\end{itemize}

No suture removal is necessary as the LASIK flap self-seals and adheres naturally to the stromal bed.

\section*{Outcomes \& Prognosis}
LASIK is one of the most frequently performed elective surgical procedures worldwide, with millions of procedures conducted annually across various continents. The incidence of LASIK has seen significant growth since its introduction, peaking in the early 2000s and remaining a popular choice for refractive correction due to its high success rates and rapid recovery. The procedure is most commonly performed on young to middle-aged adults, typically between 20 and 45 years old, who have stable refractive errors and are seeking to reduce their dependence on corrective eyewear. While there is no significant sex predilection, myopia is the most common indication, accounting for the vast majority of cases, followed by astigmatism and hyperopia. Mortality directly attributable to LASIK is virtually zero, making it an exceptionally safe procedure in terms of life-threatening risks. Morbidity, while low, primarily involves complications such as dry eye syndrome (affecting a significant percentage of patients temporarily, with a smaller subset experiencing chronic symptoms), flap-related issues (less than 1\% with modern femtosecond techniques), and the rare but serious risk of corneal ectasia (estimated at 0.02\% to 0.6\% in various cohorts).

The outcomes and prognosis for LASIK are generally excellent, with high patient satisfaction rates consistently reported in large-scale studies and meta-analyses. Studies consistently show that over 95\% of patients achieve a driving vision of 20/40 or better without correction, and approximately 85-90\% achieve 20/20 uncorrected visual acuity. Functional outcomes include a significant improvement in quality of life, increased freedom from glasses and contact lenses, and enhanced participation in sports and daily activities. The long-term stability of LASIK results is generally good, with most patients maintaining their corrected vision for many years. However, some degree of regression, where the refractive error partially returns, can occur in a small percentage of patients, particularly those with very high initial corrections or underlying ocular surface issues.

Prognostic factors influencing outcomes include the magnitude of the preoperative refractive error, corneal thickness and topography, age, and the presence of any underlying ocular conditions. While the risk of severe vision-threatening complications is very low, patients should be aware of potential long-term issues such as chronic dry eye or the rare development of corneal ectasia, which may require further intervention like corneal collagen cross-linking or corneal transplantation. The safety and efficacy of LASIK are well-documented by regulatory bodies like the FDA, which continuously monitors post-market data and clinical trials to ensure patient safety and optimal outcomes.

\section*{References}
\begin{enumerate}
  \item MedlinePlus. LASIK Eye Surgery. U.S. National Library of Medicine; 2024. Available from: https://medlineplus.gov/lasikeyesurgery.html
  
  \item Yanoff M, Duker JS. Ophthalmology. 6th ed. Elsevier; 2023. Chapter 6.10: Refractive Surgery.
  
  \item American Academy of Ophthalmology. Preferred Practice Pattern: Refractive Errors and Refractive Surgery. San Francisco, CA: American Academy of Ophthalmology; 2021.
  
  \item Azar DT, Steinert RF. LASIK: Surgical Techniques and Complications. 3rd ed. CRC Press; 2020.
\end{enumerate}

\clearpage
